% !TeX program = xelatex
% Compile: xelatex thuyetminh_de_tai.tex && xelatex thuyetminh_de_tai.tex
% Mau: THUYET MINH DE TAI NGHIEN CUU KHOA HOC SINH VIEN (theo cau truc file
% _NCKH_THUYETMINH_DETAI cua Truong Dai hoc Cong nghe Thong tin, DHQG TP.HCM)
\documentclass[a4paper,12pt]{article}

\usepackage{fontspec}
\setmainfont{Noto Serif}
\setsansfont{Noto Sans}
\usepackage{polyglossia}
\setmainlanguage{vietnamese}

\usepackage[a4paper,margin=2.3cm]{geometry}
\usepackage{amsmath,amssymb}
\usepackage{graphicx}
\usepackage{booktabs}
\usepackage{longtable}
\usepackage{array}
\usepackage{enumitem}
\usepackage{xcolor}
\usepackage{hyperref}
\hypersetup{
    colorlinks=true,
    linkcolor=blue!45!black,
    urlcolor=blue!45!black,
    citecolor=blue!45!black,
    pdftitle={Thuyet minh de tai NCKH -- SADGS Digital Twin},
    pdfauthor={Digital Twin GS}
}
\usepackage{titlesec}

\definecolor{hdblue}{RGB}{31,73,125}
\titleformat{\section}{\normalfont\large\bfseries\color{hdblue}}{\thesection.}{0.5em}{}
\titleformat{\subsection}{\normalfont\normalsize\bfseries\color{hdblue}}{\thesubsection}{0.5em}{}

\setlength{\parindent}{0pt}
\setlength{\parskip}{5pt}

\begin{document}

% ---------------------------------------------------------------------
% Trang bia / header hanh chinh
% ---------------------------------------------------------------------
\begin{minipage}[t]{0.68\textwidth}
\centering
\textsc{ĐẠI HỌC QUỐC GIA TP.HCM}\\
\textbf{TRƯỜNG ĐẠI HỌC QUỐC TẾ}
\end{minipage}
\hfill
\begin{minipage}[t]{0.28\textwidth}
\small
\begin{tabular}{|p{3.4cm}|}
\hline
Ngày nhận hồ sơ \\
\hline
\rule{0pt}{1.1cm} \\
\hline
\multicolumn{1}{c}{\footnotesize\textit{(Do CQ quản lý ghi)}} \\
\end{tabular}
\end{minipage}

\vspace{0.6cm}
\begin{center}
{\LARGE\bfseries THUYẾT MINH}\\[2pt]
{\Large\bfseries ĐỀ TÀI NGHIÊN CỨU KHOA HỌC SINH VIÊN}
\end{center}

\vspace{0.4cm}
\section*{A. THÔNG TIN CHUNG}
\vspace{-0.3cm}

\subsection*{A1. Tên đề tài}
\begin{itemize}[leftmargin=1.2em]
\item Tên tiếng Việt (IN HOA):\\
\textbf{XÂY DỰNG TÀI LIỆU DIGITAL TWIN VÀ KIỂM CHỨNG THỰC NGHIỆM CƠ CHẾ PHÂN BỔ MẬT ĐỘ NHẬN THỨC CẤU TRÚC (STRUCTURE-AWARE DENSIFICATION) TRONG 3D GAUSSIAN SPLATTING}
\item Tên tiếng Anh (IN HOA):\\
\textbf{A VERIFIED DIGITAL-TWIN DOCUMENTATION AND BENCHMARKING FRAMEWORK FOR STRUCTURE-AWARE DENSIFICATION IN 3D GAUSSIAN SPLATTING (SADGS)}
\end{itemize}

\subsection*{A2. Thời gian thực hiện}
6 tháng (kể từ khi được duyệt).

\subsection*{A3. Tổng kinh phí}
Tổng kinh phí: \textit{[điền theo mức tài trợ thực tế]}, gồm:
\begin{itemize}[leftmargin=1.2em]
\item Kinh phí từ \textit{[đơn vị tài trợ]}: \textit{[số tiền]}
\end{itemize}

\subsection*{A4. Chủ nhiệm}
\begin{tabular}{@{}ll@{}}
Họ và tên: & Lê Anh Kiệt \\
Ngày, tháng, năm sinh: & \textit{[điền]} \hspace{1.5cm} Giới tính (Nam/Nữ): \textit{[điền]} \\
Số CCCD: & \textit{[điền]} \quad ; Ngày cấp: \textit{[điền]} ; Nơi cấp: \textit{[điền]} \\
Mã số sinh viên: & ITCSIU24047 \\
Số điện thoại liên lạc: & \textit{[điền]} \\
Trường: & Đại học Quốc tế -- Đại học Quốc gia TP.HCM \\
Khoa: & \textit{[điền]} \\
Số tài khoản: & \textit{[điền]} \hspace{1.5cm} Ngân hàng: \textit{[điền]} \\
\end{tabular}

\subsection*{A5. Thành viên đề tài (kể cả CNĐT)}
\begin{longtable}{@{}p{0.6cm} p{4.2cm} p{3cm} p{5cm}@{}}
\toprule
\textbf{TT} & \textbf{Họ tên} & \textbf{MSSV} & \textbf{Trường / Khoa} \\
\midrule
\endhead
1 & Lê Anh Kiệt & ITCSIU24047 & Đại học Quốc tế -- ĐHQG TP.HCM \\
\bottomrule
\end{longtable}

\newpage
\section*{B. MÔ TẢ NGHIÊN CỨU}

\subsection*{B1. Giới thiệu về đề tài}

\textbf{Giới thiệu bài toán:} 3D Gaussian Splatting (3DGS) biểu diễn một cảnh 3D bằng một tập
hợp Gaussian dị hướng (vị trí, hiệp phương sai, độ mờ đục, hệ số màu theo bậc cầu hài -- spherical
harmonics), khởi tạo từ point cloud SfM (COLMAP) và tối ưu bằng render khả vi (differentiable
rasterization) so khớp với ảnh chụp thật. 3DGS đạt chất lượng tương đương NeRF nhưng render
thời gian thực, nên có nhiều ứng dụng:
\begin{itemize}[leftmargin=1.2em]
\item \textbf{Digital twin \& thực tại ảo/tăng cường (VR/AR):} tái tạo nhanh không gian thật
thành mô hình 3D tương tác được theo thời gian thực.
\item \textbf{Robot học \& SLAM thị giác:} biểu diễn bản đồ 3D dày đặc, khả vi, dùng cho định
vị và lập kế hoạch đường đi.
\item \textbf{Số hoá di sản văn hoá \& khảo sát công trình:} tái tạo hiện trạng 3D độ phân giải
cao từ ảnh chụp thông thường, chi phí thấp hơn quét LiDAR.
\end{itemize}

Một khâu quan trọng của pipeline 3DGS là \emph{Adaptive Density Control} (ADC): quyết định khi
nào một Gaussian cần được nhân đôi (clone), tách (split), hay loại bỏ (prune) trong quá trình
huấn luyện, dựa trên gradient vị trí tích luỹ qua các bước. Đây là yếu tố quyết định cả chất
lượng tái tạo lẫn số lượng Gaussian cuối cùng (ảnh hưởng trực tiếp tới bộ nhớ và tốc độ render).

\textbf{Lý do chọn đề tài:} ADC nguyên bản chỉ dùng độ lớn gradient vị trí trung bình -- một tín
hiệu gián tiếp, dễ bị nhiễu bởi hiện tượng \emph{gradient cancellation} khi nhiều góc nhìn (view)
``triệt tiêu'' lẫn nhau, và tách Gaussian theo kiểu đẳng hướng (isotropic) dù cấu trúc ảnh xung
quanh có thể dị hướng rõ rệt (cạnh, đường thẳng, viền vật thể). Hướng nghiên cứu gần đây --
\emph{Structure-Aware Densification for Gaussian Splatting} (SADGS, Lyu et al., SIGGRAPH 2026)
-- đề xuất thay tín hiệu gradient bằng một tín hiệu trực tiếp hơn: so khớp \emph{screen-space
extent} của mỗi Gaussian với \emph{cấu trúc texture đa tỉ lệ} (multiscale structure tensor) của
ảnh mục tiêu, suy ra tỉ số $\eta$, từ đó thực hiện \emph{split dị hướng} (anisotropic split) theo
đúng hướng cấu trúc ảnh, cộng thêm bước \emph{multiview-consistency} để tránh quyết định dựa
trên một view nhiễu đơn lẻ.

Tuy nhiên, khi khảo sát trực tiếp mã nguồn công khai của SADGS đối chiếu với tài liệu/bài báo đi
kèm, nhóm nghiên cứu nhận thấy một khoảng cách đáng kể giữa những gì \emph{được mô tả} và những
gì \emph{thực sự chạy} trong pipeline huấn luyện mặc định: một số hàm được cài đặt đầy đủ trong
mã nguồn (ví dụ \texttt{final\_prune\_structgs}, \texttt{expand\_undersized\_gs}) nhưng lời gọi
của chúng lại bị vô hiệu hoá (comment out) trong \texttt{train.py}; một số tín hiệu ``importance''
được mô tả như điểm số chất lượng render đa view nhưng thực chất chỉ là bộ đếm số lần nhìn thấy
(\texttt{accum\_view\_count}). Đây chính là động lực cho hướng tiếp cận \textbf{digital twin}
của đề tài: thay vì chỉ diễn giải lại bài báo, mọi công thức/tham số/tên hàm trình bày trong tài
liệu phải được \textbf{đối chiếu trực tiếp (grep/read) với mã nguồn thật}, phân biệt rõ ràng cơ
chế đang thực sự chạy (\emph{live path}) với phần đã cài đặt nhưng không được kích hoạt
(\emph{dead code}) -- và nội dung không truy vết được tới code/dữ liệu thật sẽ bị loại bỏ thay vì
giữ lại như một giả định chưa kiểm chứng.

\subsection*{B2. Mục tiêu, nội dung, kế hoạch nghiên cứu}

\subsubsection*{B2.1 Mục tiêu}
\begin{itemize}[leftmargin=1.2em]
\item Khảo sát pipeline 3D Gaussian Splatting gốc (khởi tạo, rasterizer khả vi, loss, backprop,
Adaptive Density Control) làm nền tảng đối chiếu.
\item Nghiên cứu và \textbf{xác minh trực tiếp với mã nguồn} cơ chế Structure-Aware Densification
của SADGS: structure tensor đa tỉ lệ, tỉ số $\eta$, split dị hướng, multiview-consistency pruning
-- xác định rõ phần nào đang chạy thật trong \texttt{train.py} và phần nào là dead code.
\item Xây dựng bộ tài liệu digital twin (sách kỹ thuật + slide + hình minh hoạ) mà mọi hình vẽ
được sinh từ dữ liệu/ảnh/checkpoint thật thay vì số liệu giả định, và mọi khẳng định kỹ thuật đều
truy vết được tới file:dòng cụ thể trong mã nguồn.
\item Thực nghiệm, chạy lại benchmark chính thức cho cơ chế densify/prune thật của SADGS (khác
với số liệu tham khảo hiện có của pipeline trước khi tích hợp), so sánh định lượng với 3DGS gốc
trên cùng bộ dữ liệu, dùng các độ đo PSNR, SSIM, LPIPS, số Gaussian cuối, thời gian huấn luyện,
VRAM đỉnh.
\end{itemize}

\subsubsection*{B2.2 Nội dung và phương pháp nghiên cứu}

\textbf{Nội dung 1: Khảo sát pipeline 3D Gaussian Splatting gốc và cơ chế Adaptive Density
Control}

\textit{Phương pháp thực hiện:}
\begin{itemize}[leftmargin=1.2em]
\item Đọc và tái hiện (bằng Python/NumPy, đối chiếu với mã CUDA thật) các khâu: khởi tạo Gaussian
từ SfM point cloud, rasterizer khả vi (tile-based, alpha blending), hàm loss ($L_1$ + D-SSIM),
lan truyền ngược, và ADC nguyên bản (\texttt{densify\_and\_clone}, \texttt{densify\_and\_split}).
\item Xác định hạn chế cụ thể của ADC gốc: gradient cancellation đa view, split đẳng hướng không
tính tới cấu trúc ảnh cục bộ.
\end{itemize}

\textit{Kết quả dự kiến:} tài liệu nền tảng toán học và kiến trúc 3DGS, có kiểm chứng công thức
qua mã nguồn, làm cơ sở đối chiếu cho SADGS ở Nội dung 2.

\medskip
\textbf{Nội dung 2: Nghiên cứu, xác minh và tài liệu hoá cơ chế Structure-Aware Densification
(SADGS)}

\textit{Phương pháp thực hiện:}
\begin{itemize}[leftmargin=1.2em]
\item Đọc mã nguồn SADGS (\texttt{scene/gaussian\_model.py}, \texttt{utils/freq\_utils.py},
\texttt{train.py}) để xác định chính xác: công thức structure tensor đa tỉ lệ (Laplacian scale
space, Di Zenzo), cách tính tỉ số $\eta = \lambda_{\max}/\lambda_{\min}$, quy tắc split dị hướng
theo trục ($k_{\text{axis}} = \lceil\sqrt{\eta_{\text{axis}}}\rceil$), và cơ chế
multiview-consistency.
\item Với mỗi cơ chế, xác minh bằng \texttt{grep}/đọc trực tiếp \texttt{train.py} xem lời gọi có
thực sự nằm trên đường thực thi mặc định hay không (ví dụ: \texttt{final\_prune\_structgs} và
\texttt{expand\_undersized\_gs} bị comment out; \texttt{importance\_score} truyền vào thực chất
là \texttt{accum\_view\_count}, không phải điểm lỗi pixel).
\item Sinh hình minh hoạ từ dữ liệu thật: tái hiện công thức bằng NumPy thuần, chạy trên ảnh huấn
luyện thật và checkpoint Gaussian thật (không dùng số liệu ngẫu nhiên/giả định để minh hoạ một
phép đo).
\item Biên soạn thành sách kỹ thuật (dạng chương/mục có đánh số) và hai bộ slide LaTeX/Beamer
(bản đầy đủ tra cứu và bản báo cáo rút gọn), mỗi khẳng định kỹ thuật kèm trích dẫn
\texttt{file:dòng} trong mã nguồn.
\end{itemize}

\textit{Kết quả dự kiến:} bộ tài liệu digital twin hoàn chỉnh cho cơ chế SADGS, phân biệt rõ cơ
chế đang chạy thật và phần dead code, làm cơ sở cho thực nghiệm định lượng ở Nội dung 3.

\medskip
\textbf{Nội dung 3: Thực nghiệm và đánh giá định lượng SADGS so với 3DGS gốc}

\textit{Phương pháp thực hiện:}
\begin{itemize}[leftmargin=1.2em]
\item Chạy huấn luyện đầy đủ (30\,000 iteration) cho cả hai cấu hình -- ADC 3DGS gốc và
Structure-Aware Densification của SADGS -- trên cùng (các) scene chuẩn, cùng hạ tầng tính toán.
\item Ghi lại số liệu theo checkpoint: PSNR, SSIM, LPIPS, số Gaussian, thời gian huấn luyện, VRAM
đỉnh (tương tự định dạng \texttt{history.csv}/\texttt{leaderboard.csv} đã dùng cho pipeline hiện
tại).
\item So sánh định lượng (bảng số liệu) và định tính (ảnh render cạnh nhau) giữa hai cấu hình;
phân tích ảnh hưởng của densify dị hướng lên vùng biên/chi tiết nhỏ.
\end{itemize}

\textit{Kết quả dự kiến:} bảng số liệu benchmark SADGS \emph{chính thức} (thay thế số liệu tham
khảo hiện có của pipeline trước khi tích hợp cơ chế densify/prune thật), cùng phân tích định
lượng/định tính về hiệu quả của structure-aware densification.

\subsubsection*{B2.3 Kế hoạch nghiên cứu}

\begin{longtable}{@{}p{1.3cm} p{3.4cm} p{9.3cm}@{}}
\toprule
\textbf{Nội dung} & \textbf{Thời gian dự kiến} & \textbf{Kết quả mong muốn} \\
\midrule
\endhead
1 & Tháng 1 -- 2 & Tài liệu khảo sát pipeline 3DGS gốc và ADC, có kiểm chứng công thức qua mã
nguồn. \\
2 & Tháng 3 -- 4 & Bộ tài liệu digital twin cho SADGS (sách + 2 bộ slide + hình minh hoạ từ dữ
liệu thật), mọi khẳng định kỹ thuật kèm trích dẫn file:dòng. \\
3 & Tháng 5 -- 6 & Bảng số liệu benchmark SADGS chính thức so với 3DGS gốc, kèm phân tích định
lượng/định tính. \\
\bottomrule
\end{longtable}

\subsection*{B3. Kết quả dự kiến}
\begin{itemize}[leftmargin=1.2em]
\item Bộ tài liệu digital twin hoàn chỉnh cho SADGS: sách kỹ thuật (dạng chương/mục), hai bộ
slide LaTeX/Beamer, và tập hình minh hoạ được sinh từ dữ liệu/ảnh/checkpoint thật -- mọi công
thức, tên hàm, tham số mặc định đều kiểm chứng được trực tiếp với mã nguồn.
\item Báo cáo so sánh định lượng và định tính giữa 3DGS gốc và SADGS (PSNR, SSIM, LPIPS, số
Gaussian, thời gian huấn luyện, VRAM), dựa trên một lượt chạy benchmark thật cho cơ chế
densify/prune của SADGS.
\item Quy trình/checklist xác minh tài liệu đối chiếu mã nguồn, có thể tái sử dụng cho các dự án
tài liệu hoá phần mềm nghiên cứu khác.
\end{itemize}

\section*{Tài liệu tham khảo}
\begin{enumerate}[leftmargin=1.4em]
\item Kerbl, B., Kopanas, G., Leimkühler, T., \& Drettakis, G. (2023). 3D Gaussian Splatting for
Real-Time Radiance Field Rendering. \textit{ACM Transactions on Graphics}, 42(4).
\item Lyu, L., Tewari, A., Chen, J., Leimkühler, T., \& Theobalt, C. (2026). Faster 3D Gaussian
Splatting Convergence via Structure-Aware Densification. \textit{SIGGRAPH 2026 Conference
Track}. \url{https://vcai.mpi-inf.mpg.de/projects/SAD-GS/}
\item FastGS -- dẫn nguồn theo \texttt{SADGS/README.md} (mục Acknowledgements) như công trình
liên quan mà SADGS kế thừa một phần tham số cấu hình; \textbf{không} phải mã nguồn được triển
khai trong repo SADGS đang khảo sát.
\end{enumerate}

\vspace{1.2cm}
\begin{minipage}{0.48\textwidth}
\centering
\textit{Ngày \phantom{00} tháng \phantom{00} năm \phantom{0000}}\\
\textbf{Giảng viên hướng dẫn}\\
\textit{(Ký và ghi rõ họ tên)}
\vspace{1.6cm}
\end{minipage}
\hfill
\begin{minipage}{0.48\textwidth}
\centering
\textit{Ngày \phantom{00} tháng \phantom{00} năm \phantom{0000}}\\
\textbf{Chủ nhiệm đề tài}\\
\textit{(Ký và ghi rõ họ tên)}
\vspace{1.6cm}
\end{minipage}

\end{document}
